\subsection{Index and Exponent}

THEOREM 6. --- Let K be a commutative field, and let A be a central simple A-algebra of finite degree over K. Let L be a separable extension of K of finite degree that is a splitting field for the algebra A. Then \([L:K][A]\) is zero in \(\mathrm{Br}(K)\) .

There exists an extension M of L that is a Galois extension of K of finite degree (V, §10, No.~1, p.~57, Proposition 2). The class {[}A{]} of A in the Brauer group of K belongs to the subgroup Br(M/K). Let G be the Galois group of M over K, and let \(\alpha\) be the image of {[}A{]} in \(H^{2}(G,M^{*})\) (VIII, p.~321, Corollary). Let H be the Galois group of M over L. Then H is a subgroup of index {[}L:K{]} in G (V, §10, No.~7, p.~68, Corollary 5). Since \(\operatorname{Res}_{\mathrm{H}}^{\mathrm{G}}(\alpha)=\Phi_{\mathrm{M/L}}(\mathrm{A}_{(\mathrm{L})})\) (Proposition 14), we have \(\operatorname{Res}_{\mathrm{H}}^{\mathrm{G}}(\alpha)=0\) . By Proposition 8 of VIII, p.~303, it follows that {[}L:K{]} \(\alpha=0\) , and consequently {[}L:K{]}{[}A{]}=0.

Let K be a commutative field, and let A be a central simple algebra of finite degree over K. Then A is isomorphic to an algebra of the form \(\mathbf{M}_{n}(\mathrm{D})\) , where D is a field with center K, and \([A]=[D]\) in \(\mathrm{Br}(K)\) . The reduced degree of D depends only on A. We call this reduced degree the index of A. The index of A divides the reduced degree of A. The exponent of A is the order of the class of A in the Brauer group of K.

COROLLARY 1. --- The exponent of a central simple algebra of finite degree over a commutative field divides its index.

It suffices to prove this for a field D of finite degree over its center K. Let L be a maximal commutative subfield of D that is a separable extension of K; we have \([D:K]=[L:K]^{2}\) (VIII, p.~265, Corollary 2, b) and c)). Then the extension L of K is a splitting field for the algebra D (VIII, p.~281, Proposition 5), and \([L:K]\) coincides with the reduced degree of D. We then apply Theorem 6.

COROLLARY 2. --- Let K be a commutative field, and let A be a central simple algebra of finite degree over K. Let p be a prime number. If p divides the index of A, then it divides the exponent of A.

Let us suppose that the prime number \(p\) does not divide the exponent of A and prove that it does not divide its index. It suffices to prove this result in the case when A is a field. Let L be a Galois extension of K of finite degree that splits A. Let G be the Galois group of L over K, and let H be a Sylow \(p\) -subgroup of G (I, §6, No.~6, p.~78). We denote by \(\mathrm{L}' = \mathrm{L}^{\mathrm{H}}\) the subfield of L of invariants of H. The exponent of \(\mathrm{A}_{(\mathrm{L}')}\) divides that of A and is therefore prime to \(p\) . By Theorem 6, we have \([\mathrm{L}:\mathrm{L}'][\mathrm{A}_{(\mathrm{L}')}] = 0\) in the Brauer group of \(\mathrm{L}'\) . It follows that \([\mathrm{A}_{(\mathrm{L}')}] = 0\) and that the field \(\mathrm{L}'\) splits A. We then apply Corollary 2 of VIII, p.~283; it follows that the index of A divides \([\mathrm{L}':\mathrm{K}]\) and is therefore not divisible by \(p\) .

COROLLARY 3. --- Let p be a prime number and let K be a perfect field of characteristic p.~Let A be a central simple algebra of finite degree over K. Then p does not divide the index of A.

Let us prove that the Brauer group of K does not contain any element of order p.~Every Galois extension M of K of finite degree is a perfect field (V, §7, No.~1, p.~36, Proposition 2). Hence taking the p-th power is an automorphism of the group \(M^{*}\) . It follows that multiplication by p is an automorphism of the group \(\mathrm{H}^{2}(\mathrm{Gal}(\mathrm{M}/\mathrm{K}),\mathrm{M}^{*})\) that is isomorphic to \(\mathrm{Br}(\mathrm{M}/\mathrm{K})\) .

Consequently, the order of {[}A{]} is prime to \(p\) and, by Corollary 2, its index is not divisible by \(p\) .

Remark. --- By considering tensor products of quaternion algebras, it is possible to construct fields with center K, exponent 2, and arbitrarily large index (cf.~VIII, p.~371, Exercises 7 and 8).

*Conversely, if K is a finite Galois extension of a \(p\) -adic field or a field of formal power series over a finite field, then the exponent of a central simple algebra of finite degree over K is equal to its index (VIII, p.~332, Exercise 17, e)).*

